\section{Preliminaries for physical renormalization}

We first collect the transfer-map, channel, and physical-closure notions used
to state the mixed-state renormalization condition.  The renormalization fixed
point of \cite[Definition~4.1]{Cirac2016MPDO_arXiv} is stated in the next
section.

\begin{definition}[MPO transfer-map idempotence]\label{def:is_rfp_mpdo}
    \lean{MPOTensor.IsRFP}
    \leanok
    \uses{def:mpo_transfer_map}
    An MPO tensor $M$ has \emph{transfer-map idempotence} if
    \[
        \E_M \circ \E_M = \E_M.
    \]
\end{definition}

% The definition above is transfer-map idempotence, distinct from the paper's
% Definition 4.1 (`RFPMixedTS`); the latter is stated below as def:is_rfp_via_ts.

\begin{definition}[Trace-preserving completely positive map]
    \label{def:is_kraus_cptp}
    \lean{IsKrausCPTP}
    \leanok
    A linear map $\mathcal{S}$ between matrix algebras is \emph{trace-preserving
        completely positive} if it has a Kraus form
    $\mathcal{S}(X) = \sum_{i} A_i X A_i^\dagger$ with $\sum_i A_i^\dagger A_i = I$.
    The Kraus operators may be rectangular, so the input and output dimensions
    need not agree.
\end{definition}

\begin{theorem}[Trace preservation from a Kraus resolution]
    \label{thm:is_kraus_cptp_trace_map}
    \lean{IsKrausCPTP.trace_map}
    \leanok
    \uses{def:is_kraus_cptp}
    Every trace-preserving completely positive map in rectangular Kraus form
    preserves the matrix trace.
\end{theorem}

\begin{proof}
    \leanok
    If $\mathcal S(X)=\sum_a A_a X A_a^\dagger$ and
    $\sum_a A_a^\dagger A_a=I$, cyclicity gives
    \[
        \tr\mathcal S(X)
        =\sum_a\tr(A_aXA_a^\dagger)
        =\tr\left(\left(\sum_a A_a^\dagger A_a\right)X\right)
        =\tr X.
    \]
\end{proof}

\begin{lemma}[Positivity preservation for Kraus maps]
    \label{lem:is_kraus_cptp_map_pos}
    \lean{IsKrausCPTP.map_posSemidef}
    \leanok
    \uses{def:is_kraus_cptp}
    If $\mathcal S$ is trace-preserving completely positive and $X\geq0$, then
    $\mathcal S(X)\geq0$.
\end{lemma}

\begin{lemma}[Identity map is trace-preserving completely positive]
    \label{lem:is_kraus_cptp_id}
    \lean{isKrausCPTP_id}
    \leanok
    \uses{def:is_kraus_cptp}
    The identity map $\mathrm{id}(X) = X$ is trace-preserving completely positive,
    with single Kraus operator $A_0 = I$.
\end{lemma}
\begin{proof}
    \leanok
    Take $r = 1$, $A_0 = I$; then $\mathrm{id}(X) = I X I^\dagger = X$ and
    $A_0^\dagger A_0 = I$.
\end{proof}

\begin{theorem}[Conjugation by an isometry is trace-preserving completely positive]
    \label{thm:mpdo_single_kraus_isometry_cptp}
    \lean{isKrausCPTP_of_singleKraus}
    \leanok
    \uses{def:is_kraus_cptp}
    Let $H$ and $K$ be finite-dimensional complex vector spaces, and let
    $V:H\to K$ satisfy $V^\dagger V=I_H$. Then the map
    \[
        \Phi_V:\operatorname{End}_{\C}(H)\longrightarrow
          \operatorname{End}_{\C}(K),
        \qquad \Phi_V(X)=VXV^\dagger,
    \]
    is trace-preserving and completely positive. This is the general
    one-isometry form of the local basis change used in
    \cite[Appendix~C.2, lines~1439 and~1520]{Cirac2016MPDO_arXiv}.
\end{theorem}
\begin{proof}
    \leanok
    Take the sole Kraus operator to be $V$. Its resolution of the identity is
    precisely $V^\dagger V=I_H$.
\end{proof}

\begin{definition}[Single-Kraus map]
    \label{def:mpdo_single_kraus_map}
    \lean{sandwichMap}
    \leanok
    For a rectangular matrix $V:H\to K$, define the linear map
    \[
      \Phi_V(X)=VXV^\dagger.
    \]
\end{definition}

\begin{lemma}[Evaluation of the single-Kraus map]
    \label{lem:mpdo_single_kraus_map_apply}
    \lean{sandwichMap_apply}
    \leanok
    \uses{def:mpdo_single_kraus_map}
    For every matrix $X$, one has $\Phi_V(X)=VXV^\dagger$.
\end{lemma}

\begin{theorem}[The single-Kraus map of an isometry is a channel]
    \label{thm:mpdo_single_kraus_map_cptp}
    \lean{sandwichMap_isKrausCPTP}
    \leanok
    \uses{def:mpdo_single_kraus_map, thm:mpdo_single_kraus_isometry_cptp}
    If $V^\dagger V=I_H$, then $\Phi_V$ is trace-preserving and completely
    positive.
\end{theorem}

\begin{definition}[Matrix reindexing along an equivalence]
    \label{def:mpdo_equiv_reindex_map}
    \lean{Matrix.equivReindexMap}
    \leanok
    Let $e:I\simeq J$ be an equivalence of index sets. The associated
    matrix reindexing is the linear map
    $R_e:\C^{I\times I}\to\C^{J\times J}$ determined by
    \[
        \bigl[R_e(X)\bigr]_{j,j'}=X_{e^{-1}(j),e^{-1}(j')}.
    \]
\end{definition}

\begin{theorem}[Equivalence reindexing is trace-preserving completely positive]
    \label{thm:mpdo_equiv_reindex_cptp}
    \lean{Matrix.equivReindexMap_isKrausCPTP}
    \leanok
    \uses{def:mpdo_equiv_reindex_map, thm:mpdo_single_kraus_isometry_cptp}
    For every equivalence $e:I\simeq J$ between finite index sets, the
    reindexing map $R_e$ is trace-preserving and completely positive.
\end{theorem}
\begin{proof}
    \leanok
    Let $P_e:\C^I\to\C^J$ be the permutation matrix of $e$. Then
    \[
        R_e(X)=P_eXP_e^\dagger,
        \qquad P_e^\dagger P_e=I_{\C^I}.
    \]
    The claim follows from the preceding isometry lemma.
\end{proof}

% The concrete dependent-sum equivalences for T_0, S_0, T_2, and S_2 are
% stated in ch21_mpdo_rfp_simple_local_normalized_preparations.tex after the eta preparations.

\begin{lemma}[Composition of trace-preserving completely positive maps]
    \label{lem:is_kraus_cptp_comp}
    \lean{isKrausCPTP_comp}
    \leanok
    \uses{def:is_kraus_cptp}
    The composition of two trace-preserving completely positive maps is again
    trace-preserving completely positive. If $\mathcal{S}$ has Kraus operators
    $A_i$ and $\mathcal{T}$ has Kraus operators $B_j$, then $\mathcal{S} \circ
    \mathcal{T}$ has Kraus operators $A_i B_j$.
\end{lemma}
\begin{proof}
    \leanok
    The Kraus form of the composite is
    $(\mathcal{S} \circ \mathcal{T})(X) = \sum_{i,j} (A_i B_j) X (A_i B_j)^\dagger$,
    and the resolution of identity follows from those of $\mathcal{S}$ and
    $\mathcal{T}$:
    \[
        \sum_{i,j} (A_i B_j)^\dagger (A_i B_j)
        = \sum_{j} B_j^\dagger \Bigl( \sum_{i} A_i^\dagger A_i \Bigr) B_j
        = \sum_{j} B_j^\dagger B_j = I.
    \]
\end{proof}

\begin{theorem}[Tensor extension preserves trace-preserving complete positivity]
    \label{thm:mpdo_tensor_map_id_cptp}
    \lean{Matrix.tensorMapIdLM_isKrausCPTP}
    \leanok
    \uses{def:tensor_map_id, def:is_kraus_cptp}
    Let $\mathcal S:\operatorname{End}_{\C}(H)\to
        \operatorname{End}_{\C}(K)$ be trace-preserving and completely positive,
    and let $R$ be another finite-dimensional space. Then
    \[
        \mathcal S\otimes\id_R:
        \operatorname{End}_{\C}(H\otimes R)\longrightarrow
        \operatorname{End}_{\C}(K\otimes R)
    \]
    is trace-preserving and completely positive. If $(A_a)_a$ is a Kraus
    family for $\mathcal S$, then
    $(A_a\otimes I_R)_a$ is a Kraus family for its tensor extension.
\end{theorem}
\begin{proof}
    \leanok
    The tensor extension has the Kraus form
    \[
        (\mathcal S\otimes\id_R)(X)
        =\sum_a (A_a\otimes I_R)X(A_a\otimes I_R)^\dagger.
    \]
    Its normalization follows from that of the original family:
    \[
        \sum_a (A_a\otimes I_R)^\dagger(A_a\otimes I_R)
        =\left(\sum_a A_a^\dagger A_a\right)\otimes I_R
        =I_H\otimes I_R.
    \]
\end{proof}

\begin{definition}[Rectangular Kraus map]
    \label{def:mpdo_rectangular_kraus_map}
    \lean{Matrix.rectangularKrausMap}
    \leanok
    Let $I$ be a finite index set, let $J$ be an arbitrary index set, and let
    $(A_a\in\C^{J\times I})_a$ be a finite family of matrices. Its
    rectangular Kraus map is
    \[
        \Phi_A(X)=\sum_a A_aXA_a^{\dagger}.
    \]
\end{definition}

\begin{theorem}[Rectangular Kraus resolution]
    \label{thm:mpdo_rectangular_kraus_map_cptp}
    \lean{Matrix.rectangularKrausMap_isKrausCPTP}
    \leanok
    \uses{def:is_kraus_cptp, def:mpdo_rectangular_kraus_map}
    Let $(A_a:H\to K)_a$ be a finite family of rectangular operators. If
    \[
        \sum_a A_a^\dagger A_a=I_H,
    \]
    then $X\mapsto\sum_a A_aXA_a^\dagger$ is trace-preserving and completely
    positive.
\end{theorem}
\begin{proof}\leanok
    The displayed family is already a Kraus representation, and the assumed
    identity is precisely its trace-preserving normalization.
\end{proof}

\begin{theorem}[Relabeling a rectangular Kraus family]
    \label{thm:mpdo_rectangular_kraus_map_equiv}
    \lean{Matrix.rectangularKrausMap_equiv}
    \leanok
    \uses{def:mpdo_rectangular_kraus_map}
    If $e:I\simeq I'$ and $A'_b=A_{e^{-1}(b)}$, then
    $\Phi_{A'}=\Phi_A$.
\end{theorem}
\begin{proof}
    \leanok
    Reindexing the finite sum gives
    \[
        \sum_{b\in I'}A_{e^{-1}(b)}X A_{e^{-1}(b)}^\dagger
        =\sum_{a\in I}A_aXA_a^\dagger.
    \]
\end{proof}

\begin{definition}[Orthogonally controlled Kraus map]
    \label{def:mpdo_controlled_kraus_map}
    \lean{Matrix.controlledKrausMap}
    \lean{Matrix.controlledKrausMap_apply_of_ne}
    \leanok
    \uses{def:mpdo_rectangular_kraus_map}
    Let $H=\bigoplus_k H_k$ and $K=\bigoplus_k K_k$. For each $k$, let
    $(A_{k,a}:H_k\to K_k)_a$ be a finite family of operators, and let
    $\widetilde A_{k,a}:H\to K$ agree with $A_{k,a}$ on $H_k$ and vanish on
    every other summand. The orthogonally controlled Kraus map is
    \[
        \mathcal C(X)
        =\sum_{k,a}\widetilde A_{k,a}X\widetilde A_{k,a}^{\dagger}.
    \]
    In particular, $\mathcal C(X)_{kl}=0$ for $k\ne l$. This is the sector
    control in the definitions of $\mathcal T_1$ and $\mathcal S_1$ in
    \cite[Appendix~C.2, lines~1523--1535 and~1548--1555]
      {Cirac2016MPDO_arXiv}.
\end{definition}

\begin{theorem}[Diagonal blocks of an orthogonally controlled map]
    \label{thm:mpdo_controlled_kraus_map_same_block}
    \lean{Matrix.controlledKrausMap_sameBlock_apply}
    \leanok
    \uses{def:mpdo_controlled_kraus_map,
      def:mpdo_rectangular_kraus_map}
    Let $X_{kk}:H_k\to H_k$ denote the $k$th diagonal block of $X$. Then
    \[
        \mathcal C(X)_{kk}
        =\sum_a A_{k,a}X_{kk}A_{k,a}^\dagger
        =\Phi_{A_k}(X_{kk}).
    \]
\end{theorem}
\begin{proof}
    \leanok
    If $j\ne k$, then every zero-extended Kraus operator satisfies
    \[
      [\widetilde A_{j,a}X\widetilde A_{j,a}^\dagger]
        _{(k,b),(k,c)}=0.
    \]
    Hence only the $j=k$ terms survive, and
    \[
      [\mathcal C(X)]_{kk}
      =\sum_a A_{k,a}X_{kk}A_{k,a}^\dagger.
    \]
\end{proof}

\begin{theorem}[Off-diagonal blocks of an orthogonally controlled map]
    \label{thm:mpdo_controlled_kraus_map_off_block}
    \lean{Matrix.controlledKrausMap_apply_of_ne}
    \leanok
    \uses{def:mpdo_controlled_kraus_map}
    If $k\ne l$, then $\mathcal C(X)_{kl}=0$.
\end{theorem}
\begin{proof}
    \leanok
    For every $(j,a)$ and $k\ne l$, the zero-extended operator satisfies
    \[
      [\widetilde A_{j,a}X\widetilde A_{j,a}^\dagger]
        _{(k,b),(l,c)}=0.
    \]
    Therefore $[\mathcal C(X)]_{kl}=\sum_{j,a}0=0$.
\end{proof}

\begin{theorem}[Orthogonal control preserves trace-preserving complete positivity]
    \label{thm:mpdo_controlled_kraus_map_cptp}
    \lean{Matrix.controlledKrausMap_isKrausCPTP}
    \leanok
    \uses{def:mpdo_controlled_kraus_map,
        def:is_kraus_cptp}
    Suppose that for every $k$ the sectorwise Kraus family resolves the
    identity:
    \[
        \sum_a A_{k,a}^{\dagger}A_{k,a}=I_{H_k}.
    \]
    Then the orthogonally controlled map $\mathcal C$ is
    trace-preserving and completely positive.
\end{theorem}
\begin{proof}
    \leanok
    \uses{thm:mpdo_rectangular_kraus_map_cptp}
    Each embedded operator has support in one summand, and therefore
    \[
        \sum_{k,a}\widetilde A_{k,a}^{\dagger}\widetilde A_{k,a}
        =\bigoplus_k\left(\sum_a A_{k,a}^{\dagger}A_{k,a}\right)
        =\bigoplus_k I_{H_k}=I_H.
    \]
    Apply Theorem~\ref{thm:mpdo_rectangular_kraus_map_cptp} to the combined
    rectangular Kraus family $(\widetilde A_{k,a})_{k,a}$.
\end{proof}

\begin{definition}[Right partial trace]
    \label{def:mpdo_right_partial_trace_map}
    \lean{Matrix.partialTraceRightLM}
    \leanok
    For a matrix $X$ on $A\otimes B$, the right partial trace is the linear map
    from matrices on $A\otimes B$ to matrices on $A$ given by
    \[
        \bigl[\tr_{B}(X)\bigr]_{ij}
        =
        \sum_{k}X_{(i,k),(j,k)}.
    \]
    After the retained and discarded subspins have been regrouped as
    $A\otimes B$, this is the partial-trace ingredient of the maps
    $\mathcal T_0$ and $\mathcal S_0$ in
    \cite[Appendix~C.2, lines~1521--1522 and~1547]{Cirac2016MPDO_arXiv}.
\end{definition}

\begin{theorem}[Right partial trace of a Kronecker product]
    \label{thm:mpdo_right_partial_trace_kronecker}
    \lean{Matrix.partialTraceRight_kronecker}
    \leanok
    \uses{def:mpdo_right_partial_trace_map}
    For square matrices $X$ on $A$ and $Y$ on $B$,
    \[
        \tr_B(X\otimes Y)=\tr(Y)X.
    \]
\end{theorem}
\begin{proof}
    \leanok
    For all $i,j$,
    \[
      [\tr_B(X\otimes Y)]_{ij}
      =\sum_t X_{ij}Y_{tt}=\tr(Y)X_{ij}.
    \]
\end{proof}

\begin{lemma}[The right partial trace is trace-preserving completely positive]
    \label{lem:mpdo_right_partial_trace_cptp}
    \lean{Matrix.partialTraceRightLM_isKrausCPTP}
    \leanok
    \uses{def:mpdo_right_partial_trace_map,
        def:is_kraus_cptp}
    The map $\tr_{B}$ is trace-preserving and completely positive for arbitrary
    finite-dimensional spaces $A$ and $B$.
\end{lemma}

\begin{proof}
    \leanok
    \uses{thm:mpdo_rectangular_kraus_map_cptp}
    For each basis vector $e_k$ of $B$, define
    $V_k:A\otimes B\to A$ by
    \[
        V_k(e_i\otimes e_\ell)=\delta_{k\ell}e_i.
    \]
    Then
    \[
        \sum_k V_k X V_k^\dagger=\tr_{B}(X),
        \qquad
        \sum_k V_k^\dagger V_k=I_{A\otimes B}.
    \]
    Theorem~\ref{thm:mpdo_rectangular_kraus_map_cptp} applies.
\end{proof}

\begin{definition}[State-preparation map]
    \label{def:mpdo_state_preparation_map}
    \lean{Matrix.preparationMap}
    \leanok
    Let $\rho$ be a matrix on a finite-dimensional space $B$.  The
    state-preparation map from matrices on $A$ to matrices on $A\otimes B$ is
    \[
        \mathcal P_\rho(X)=X\otimes\rho.
    \]
    This is the elementary preparation operation used in the maps
    $\mathcal T_1$ and $\mathcal S_1$ of
    \cite[Appendix~C.2, lines~1527--1533 and~1551--1555]{Cirac2016MPDO_arXiv}.
\end{definition}

\begin{definition}[Preparation Kraus operators]
    \label{def:mpdo_state_preparation_kraus}
    \lean{Matrix.preparationKraus}
    \leanok
    Let $R=\sqrt\rho$. For an orthonormal basis $(e_j)_j$ of $B$, define
    rectangular operators $A_j:A\to A\otimes B$ by
    \[
        A_j(e_a)=e_a\otimes Re_j.
    \]
\end{definition}

\begin{theorem}[Kraus action of state preparation]
    \label{thm:mpdo_state_preparation_kraus_action}
    \lean{Matrix.rectangularKrausMap_preparationKraus_eq}
    \leanok
    \uses{def:mpdo_rectangular_kraus_map,
      def:mpdo_state_preparation_map,
      def:mpdo_state_preparation_kraus}
    If $\rho\succeq0$ and $A_j(e_a)=e_a\otimes\sqrt\rho\,e_j$, then
    \[
        \sum_j A_jXA_j^\dagger=X\otimes\rho.
    \]
\end{theorem}
\begin{proof}
    \leanok
    Put $R=\sqrt\rho$. Since $R$ is Hermitian and $R^2=\rho$, the
    $((a,s),(b,t))$ entry of the left-hand side is
    \[
        X_{ab}\sum_jR_{sj}\overline{R_{tj}}
        =X_{ab}(R^2)_{st}=X_{ab}\rho_{st}.
    \]
\end{proof}

\begin{theorem}[Preparation Kraus operators resolve the identity]
    \label{thm:mpdo_state_preparation_kraus_resolution}
    \lean{Matrix.preparationKraus_resolution}
    \leanok
    \uses{def:mpdo_state_preparation_kraus}
    If $\rho\geq0$ and $\tr(\rho)=1$, then
    \[
        \sum_j A_j^\dagger A_j=I_A.
    \]
\end{theorem}
\begin{proof}
    \leanok
    Since $R$ is Hermitian and $R^2=\rho$, the diagonal entries of the sum are
    \[
        \sum_{j,t}\overline{R_{tj}}R_{tj}
        =\tr(R^2)=\tr(\rho)=1,
    \]
    while its off-diagonal entries vanish.
\end{proof}

\begin{lemma}[State preparation is trace-preserving completely positive]
    \label{lem:mpdo_state_preparation_cptp}
    \lean{Matrix.preparationMap_isKrausCPTP}
    \leanok
    \uses{def:mpdo_state_preparation_map,
      def:mpdo_state_preparation_kraus,
      def:is_kraus_cptp}
    If $\rho\geq0$ and $\tr(\rho)=1$, then
    $\mathcal P_\rho:X\mapsto X\otimes\rho$ is trace-preserving completely
    positive.
\end{lemma}

\begin{proof}
    \leanok
    \uses{thm:mpdo_rectangular_kraus_map_cptp,
      thm:mpdo_state_preparation_kraus_action,
      thm:mpdo_state_preparation_kraus_resolution}
    Theorem~\ref{thm:mpdo_state_preparation_kraus_action} identifies the map
    with the rectangular Kraus family $(A_j)_j$, and
    Theorem~\ref{thm:mpdo_state_preparation_kraus_resolution} gives
    $\sum_j A_j^\dagger A_j=I_A$. Apply
    Theorem~\ref{thm:mpdo_rectangular_kraus_map_cptp}.
\end{proof}

\begin{definition}[Physical closure maps]
    \label{def:phys_close}
    \lean{MPOTensor.physCloseN, MPOTensor.physCloseN_apply,
      MPOTensor.physClose1,
      MPOTensor.physClose2, MPOTensor.physClose3}
    \leanok
    \uses{def:mpo_eval_word}
    For an MPO tensor $M$, a length $N$, and a virtual operator $X$, define the
    physical operator $M_N(X)$ by
    \[
        M_N(X)_{\boldsymbol i,\boldsymbol j}
        =
        \tr\!\left(
          M^{i_0j_0} M^{i_1j_1}\cdots M^{i_{N-1}j_{N-1}} X
        \right).
    \]
    This is linear in $X$. The displayed order chooses the cut of the cyclic
    virtual contraction immediately before the first site, so that $X$ follows
    the last site. The one- and two-site operators are constructed in
    \cite[lines~638--654 and Definition~4.1]{Cirac2016MPDO_arXiv}. When
    $M=\mathcal K$, write these operators as $\mathcal K_N(X)$; the cases
    $N=2,3$ occur in
    \cite[Proposition~C.7, lines~1510--1516]{Cirac2016MPDO_arXiv}.
\end{definition}

\begin{lemma}[Periodic closure]
    \label{lem:phys_close_identity}
    \lean{MPOTensor.physCloseN_identity_eq_mpo}
    \leanok
    \uses{def:phys_close, def:mpo_family}
    Closing the virtual legs with the identity gives the periodic MPO
    operator:
    \[
        M_N(1)=\rho^{(N)}(M).
    \]
\end{lemma}

\begin{proof}
    \leanok
    For every pair of physical words $\sigma,\tau$,
    \[
        M_N(1)_{\sigma,\tau}
          = \tr\!\left(M^{\sigma,\tau}\cdot 1\right)
          = \tr\!\left(M^{\sigma,\tau}\right)
          = \rho^{(N)}(M)_{\sigma,\tau}. \qedhere
    \]
\end{proof}

\begin{theorem}[One- and two-site physical closures]
    \label{thm:phys_close_one_two}
    \lean{MPOTensor.physCloseN_one_eq_physClose1,
      MPOTensor.physCloseN_two_eq_physClose2}
    \leanok
    \uses{def:phys_close}
    Under the canonical identifications of one-site configurations with
    physical indices and two-site configurations with pairs of physical
    indices,
    \[
        M_1(X)_{i,j}=\tr(M^{ij} X),
        \qquad
        M_2(X)_{(i_0,i_1),(j_0,j_1)}
        =\tr(M^{i_0j_0} M^{i_1j_1} X).
    \]
\end{theorem}
\begin{proof}
    \leanok
    Under these identifications, the length-one and length-two word
    evaluations are $M^{ij}$ and $M^{i_0j_0}M^{i_1j_1}$, respectively.
\end{proof}

\begin{lemma}[Three-site physical closure]
    \label{lem:phys_close_three}
    \lean{finThreeArrowEquiv,
      MPOTensor.physCloseN_three_apply, MPOTensor.physClose3_apply,
      MPOTensor.physCloseN_three_eq_physClose3}
    \leanok
    \uses{def:phys_close}
    The general three-site physical closure has coefficients
    \[
        M_3(X)_{(i_0,(i_1,i_2)),(j_0,(j_1,j_2))}
        =\tr(M^{i_0j_0} M^{i_1j_1} M^{i_2j_2} X).
    \]
    For $M=\mathcal K$, this is the operator $\mathcal K_3(X)$ in
    \cite[Proposition~C.7, lines~1510--1516]{Cirac2016MPDO_arXiv}.
\end{lemma}
\begin{proof}
    \leanok
    The length-three word evaluation is
    $M^{i_0j_0}M^{i_1j_1}M^{i_2j_2}$.
\end{proof}

\section{Renormalization fixed points}

The following is the mixed-state fixed-point condition of
\cite[Definition~4.1]{Cirac2016MPDO_arXiv}.  Its two local channels extend to
every longer periodic chain by acting on the first one or two sites and leaving
the remaining sites unchanged.

\begin{definition}[Renormalization fixed point via trace-preserving maps]
    \label{def:is_rfp_via_ts}
    \lean{MPOTensor.IsRFPViaTS}
    \leanok
    \uses{def:mpo_tensor, def:is_kraus_cptp, def:phys_close}
    An MPO tensor $M$ is a \emph{renormalization fixed point} if there exist two
    trace-preserving completely positive maps $\mathcal{S}, \mathcal{T}$ on the
    physical indices such that $\mathcal{S}[M_2(X)] = M_1(X)$ and
    $\mathcal{T}[M_1(X)] = M_2(X)$ for every virtual operator $X$, where
    $M_1(X)_{ij} = \tr(M^{ij} X)$ and
    $M_2(X)_{(i_1 i_2)(j_1 j_2)} = \tr(M^{i_1 j_1} M^{i_2 j_2} X)$ are the one-
    and two-site physical operators obtained by closing one or two tensors with
    $X$. This condition is distinct from idempotence of the doubled-index
    transfer map for general mixed states. Its source zero-correlation-length
    consequence concerns the physical-trace transfer and is proved below.
    This is \cite[Definition~4.1, line~657]{Cirac2016MPDO_arXiv}.
    The two maps act as follows on every virtual closure $X$:
    \TNMPORenormalizationTS
\end{definition}

\begin{definition}[Initial-coordinate regroupings]
    \label{def:finite_tuple_regroupings}
    \lean{finSuccArrowEquiv,
      finAddTwoArrowEquiv}
    \leanok
    Let $R_1$ identify an $(N+1)$-tuple with its first coordinate and the
    remaining $N$-tuple.  Let $R_2$ identify an $(N+2)$-tuple with its first
    two coordinates and the remaining $N$-tuple.
\end{definition}

\begin{definition}[Renormalization maps on a longer ring]
    \label{def:mpdo_global_renormalization_maps}
    \lean{MPOTensor.refineFirstSite,
      MPOTensor.coarsenFirstTwoSites}
    \leanok
    \uses{def:finite_tuple_regroupings}
    Let $N$ count the sites that are left unchanged.  For physical maps
    $\mathcal T:M_d\to M_{d^2}$ and $\mathcal S:M_{d^2}\to M_d$, define
    \[
       \widehat{\mathcal T}_N
       =\mathcal T\otimes\id_{d^N},
       \qquad
       \widehat{\mathcal S}_N
       =\mathcal S\otimes\id_{d^N},
    \]
    using $R_1$ and $R_2$.  Thus $N=0$ leaves no spectator site, but the
    maps still act on a one-site or two-site operator; no empty ring occurs.
\end{definition}

\begin{lemma}[Regrouping the first physical sites]
    \label{lem:mpdo_global_renormalization_regrouping}
    \lean{finSuccArrowEquiv_apply,
      finSuccArrowEquiv_symm_apply,
      finAddTwoArrowEquiv_apply,
      finAddTwoArrowEquiv_symm_apply}
    \leanok
    \uses{def:finite_tuple_regroupings}
    For a physical word $\sigma$ of length $N+1$,
    \[
      R_1(\sigma)=\bigl(\sigma_0,(\sigma_1,\ldots,\sigma_N)\bigr),
      \qquad
      R_1^{-1}\bigl(i,(a_1,\ldots,a_N)\bigr)=(i,a_1,\ldots,a_N).
    \]
    For a physical word $\sigma$ of length $N+2$,
    \[
      R_2(\sigma)
        =\bigl((\sigma_0,\sigma_1),(\sigma_2,\ldots,\sigma_{N+1})\bigr),
    \]
    and
    \[
      R_2^{-1}\bigl((i,j),(a_1,\ldots,a_N)\bigr)
        =(i,j,a_1,\ldots,a_N).
    \]
\end{lemma}

\begin{proof}
    \leanok
    All four identities hold by definition of the product-index equivalences.
\end{proof}

\begin{lemma}[Localized renormalization maps are channels]
    \label{lem:mpdo_global_renormalization_maps_cptp}
    \lean{MPOTensor.refineFirstSite_isKrausCPTP,
      MPOTensor.coarsenFirstTwoSites_isKrausCPTP}
    \leanok
    \uses{def:mpdo_global_renormalization_maps, def:is_kraus_cptp}
    If $\mathcal S$ and $\mathcal T$ are trace-preserving completely positive,
    then so are $\widehat{\mathcal S}_N$ and
    $\widehat{\mathcal T}_N$ for every $N$.
\end{lemma}

\begin{proof}
    \leanok
    If $(K_j)_j$ are Kraus operators for $\mathcal T$ with
    $\sum_j K_j^\dagger K_j=I_d$, then the localized map has Kraus operators
    $K_j\otimes I_{d^N}$, and
    \[
      \sum_j (K_j\otimes I_{d^N})^\dagger(K_j\otimes I_{d^N})
        =\left(\sum_j K_j^\dagger K_j\right)\otimes I_{d^N}
        =I_d\otimes I_{d^N}.
    \]
    The two regroupings are unitary conjugations and preserve complete
    positivity and the trace.  The argument for $\mathcal S$ is identical.
\end{proof}

\begin{theorem}[Global action of the renormalization maps]
    \label{thm:mpdo_global_renormalization_maps}
    \lean{MPOTensor.refineFirstSite_physCloseN,
      MPOTensor.coarsenFirstTwoSites_physCloseN,
      MPOTensor.refineFirstSite_mpo,
      MPOTensor.coarsenFirstTwoSites_mpo,
      MPOTensor.exists_global_renormalization_maps}
    \leanok
    \uses{def:phys_close, def:is_rfp_via_ts,
      def:mpdo_global_renormalization_maps}
    Suppose that $\mathcal S[M_2(X)]=M_1(X)$ and
    $\mathcal T[M_1(X)]=M_2(X)$ for every virtual operator $X$.  Then, for
    every $N$ and $X$,
    \[
      \widehat{\mathcal S}_N[M_{N+2}(X)]=M_{N+1}(X),
      \qquad
      \widehat{\mathcal T}_N[M_{N+1}(X)]=M_{N+2}(X).
    \]
    In particular, these identities hold for every periodic MPO operator.
    Hence the two maps supplied by a renormalization fixed point act on rings
    of every length as required in the proof of
    \cite[Appendix~C, lines~1333--1341]{Cirac2016MPDO_arXiv}.
\end{theorem}

\begin{proof}
    \leanok
    \uses{lem:mpdo_global_renormalization_regrouping,
      lem:phys_close_identity}
    For words $u,v$ on the $N$ spectator sites, write $M^{u,v}$ for their
    virtual word evaluation and write $[Y]^{(r)}_{u,v}$ for the block of an
    operator $Y$ on the first $r$ sites.  Regrouping the first site gives
    \[
      [M_{N+1}(X)]^{(1)}_{u,v}=M_1(M^{u,v}X).
    \]
    Therefore
    \[
      \mathcal T\bigl([M_{N+1}(X)]^{(1)}_{u,v}\bigr)
        =M_2(M^{u,v}X)
        =[M_{N+2}(X)]^{(2)}_{u,v}.
    \]
    Equality of all blocks proves the refinement identity.  Interchanging
    $M_1$, $M_2$ and using $\mathcal S[M_2(Y)]=M_1(Y)$ proves the coarsening
    identity.  Setting $X=1$ and using
    $M_N(1)=\rho^{(N)}(M)$
    (Lemma~\ref{lem:phys_close_identity}) gives the periodic-ring identities
    \[
      \widehat{\mathcal T}_N[\rho^{(N+1)}(M)]
        =\rho^{(N+2)}(M),\qquad
      \widehat{\mathcal S}_N[\rho^{(N+2)}(M)]
        =\rho^{(N+1)}(M).
    \]
\end{proof}

\section{Pure-state recovery inside the MPO formalism}

\begin{definition}[Diagonal pure-state embedding as an MPO]\label{def:to_mpo_tensor}
    \lean{MPSTensor.toMPOTensor}
    \leanok
    \uses{def:mps_tensor, def:mpo_tensor}
    An MPS tensor $A = \{A^i\}_{i=0}^{d-1}$ determines an MPO tensor by
    placing $A^i$ on the diagonal in the bra and ket indices:
    \[
        M^{ij} = \delta_{ij} A^i.
    \]
    Equivalently, $M^{ii} = A^i$ and $M^{ij} = 0$ for $i \ne j$.
\end{definition}

\begin{theorem}[Transfer map of the diagonal pure-state embedding]%
    \label{thm:to_mpo_transfer_map}
    \lean{MPSTensor.toMPOTensor_transferMap}
    \leanok
    \uses{def:to_mpo_tensor, def:mpo_transfer_map, def:transfer_map}
    The transfer map of the diagonal MPO associated to $A$ agrees with the
    original MPS transfer map:
    \[
        \E_{A^{\mathrm{MPO}}} = \E_A.
    \]
\end{theorem}

\begin{proof}
    \leanok
    \uses{def:to_mpo_tensor, def:mpo_transfer_map, def:transfer_map}
    In the double sum defining $\E_{A^{\mathrm{MPO}}}$, all off-diagonal terms
    vanish because $M^{ij} = 0$ for $i \ne j$, while the diagonal terms are
    exactly $A^i X (A^i)^\dagger$. Thus only the sum over $i$ remains, which
    is the defining formula for $\E_A$.
\end{proof}

\begin{theorem}[Pure-state recovery of transfer-map idempotence]%
    \label{thm:to_mpo_rfp_iff_rfp}
    \lean{MPSTensor.toMPOTensor_isRFP_iff_isTransferIdempotent}
    \leanok
    \uses{def:is_rfp_mpdo, def:mps_transfer_idempotence,
        thm:to_mpo_transfer_map}
    For an MPS tensor $A$, the diagonal MPO embedding has idempotent transfer
    map if and only if $A$ is an RFP:
    \[
        \E_{A^{\mathrm{MPO}}}\circ\E_{A^{\mathrm{MPO}}}
        = \E_{A^{\mathrm{MPO}}}
        \iff
        \E_A\circ\E_A=\E_A.
    \]
\end{theorem}

\begin{proof}
    \leanok
    \uses{def:is_rfp_mpdo, def:mps_transfer_idempotence,
        thm:to_mpo_transfer_map}
    Both conditions assert idempotence of the corresponding transfer map.
    Theorem~\ref{thm:to_mpo_transfer_map} identifies those transfer maps, so
    the two predicates are equivalent.
\end{proof}

\begin{theorem}[Pure-state recovery of zero correlation length]%
    \label{thm:to_mpo_rfp_iff_zcl}
    \lean{MPSTensor.toMPOTensor_isRFP_iff_isZCL}
    \leanok
    \uses{thm:to_mpo_rfp_iff_rfp, thm:zcl_iff_idempotent_transfer}
    For an MPS tensor $A$, the diagonal MPO embedding has idempotent transfer
    map if and only if $A$ has zero correlation length.
\end{theorem}

\begin{proof}
    \leanok
    \uses{thm:to_mpo_rfp_iff_rfp, thm:zcl_iff_idempotent_transfer}
    Combine Theorem~\ref{thm:to_mpo_rfp_iff_rfp} with the characterization of
    zero correlation length by idempotence of the MPS transfer map
    (Theorem~\ref{thm:zcl_iff_idempotent_transfer}).
\end{proof}
